\documentclass{article}
\usepackage[utf8]{inputenc}
\usepackage{amsmath}
\usepackage{amssymb}
\usepackage{parskip}
\usepackage{csquotes}
\usepackage{microtype}

\title{Collective Decision-Making as a Learning Problem}
\author{Albert Jan van Hoek}
\date{\today}

\begin{document}

\maketitle

\section*{One More Attempt}

What if collective decision-making is really a learning problem?

I have noticed that I have become unusually invested in this problem. Partly because I believe I see a solution, and partly because I seem to keep failing to articulate what I see in a way that makes it possible for other people to see it too.

So I want to try one more time.

\section*{Starting from My Profession}

This time I will start from my own profession.

I work as an infectious disease modeller and epidemiologist in a national public health institute. Ultimately, my job is about reducing disease burden in populations.

Some of the interventions we work with are remarkably effective. Measles vaccination is one example. Measles is also one of the diseases for which global eradication is conceivable.

And eradication creates a peculiar decision problem.

If we eradicate measles, we do not merely reduce disease burden for the people alive today. We change the world inherited by every future generation. Eventually, vaccination against measles would no longer be necessary either.

The important point for my argument is not whether this future benefit can be assigned some particular expected value. It is that eradication changes the future decision problem itself.

\section*{The Decision Process}

To eradicate measles, we need a long sequence of connected decisions:

\begin{itemize}
    \item We need surveillance.
    \item We need models.
    \item We need vaccination.
    \item We need to observe what happens.
    \item We need to discover where our assumptions were wrong.
    \item We need to change the strategy.
    \item We need to coordinate with other countries.
    \item We need to learn from failures elsewhere.
\end{itemize}

The decision is therefore not really:

\begin{quote}
    A or B?
\end{quote}

It is a continuous process:

\begin{quote}
    What do we currently believe? What should we do given that belief? What happened? What did we get wrong? What does this tell us about the world? What should we believe now? What should we do next?
\end{quote}

\section*{The Cluelessness Problem}

This brings me back to the cluelessness problem.

We live in an evolving universe. Our models are necessarily incomplete, and the future is not simply unknown in the sense of being a hidden answer waiting to be calculated. The world itself changes.

So perhaps the central problem of decision-making is not how to make the correct decision from a fixed model.

Perhaps it is how to remain capable of correcting the model on which the next decision will depend.

\section*{Recursive Learning}

And this introduces another layer.

We can learn about the problem.

But we can also learn about \textit{how we learn} about the problem.

For example, better surveillance may produce better information. Better information may improve our model. A better model may lead to better decisions. The results of those decisions provide new information, which can improve the surveillance and modelling process again.

So there is a recursive possibility:

\begin{quote}
    We don't only optimise the solution; we improve the process by which we discover what the solution should be.
\end{quote}

That is the point I was trying to get at in my previous post.

But there is a second part of the problem that I had not made sufficiently explicit.

\section*{The Collective Dimension}

Measles eradication is not an individual decision problem.

\begin{itemize}
    \item No individual can eradicate measles.
    \item No single institution can eradicate measles.
    \item It requires collective action.
\end{itemize}

And now the problem becomes considerably harder.

We need a collective that can continuously adapt its understanding of a changing world.

But nobody possesses the complete model.

Different people see different things. Different institutions have different information. Different models will be wrong in different ways.

So perhaps the objective should not be to make everybody hold the same model.

Perhaps the objective is to maintain a network in which different imperfect models can continuously correct one another.

\section*{Reimagining Alignment}

This changes what I think ``alignment'' means.

\begin{itemize}
    \item We cannot necessarily align people on every proposition about the world.
    \item We cannot even guarantee that we will initially agree on the correct solution.
    \item But perhaps we can align on the process by which we discover that we are wrong.
\end{itemize}

This is where I think something interesting happens.

The things required for this process sound, at first, like ordinary virtues:

\begin{itemize}
    \item Be honest.
    \item Listen.
    \item Be curious.
    \item Say when you are uncertain.
    \item Correct others when you believe they are wrong.
    \item Allow yourself to be corrected.
    \item Distinguish what you observed from what you inferred.
    \item Help others understand information they do not have.
    \item Forgive mistakes sufficiently that people remain willing to participate.
    \item Remain connected even when you disagree.
\end{itemize}

But viewed from the perspective of a distributed learning system, these are not merely virtues.

They are \textit{functional properties} of the system.

\begin{itemize}
    \item Honesty protects the fidelity of the information entering the network.
    \item Listening determines whether information actually reaches another person's model.
    \item Correction provides an error signal.
    \item Curiosity drives exploration of uncertainty.
    \item Forgiveness helps preserve the connections through which future information can travel.
    \item Willingness to be corrected keeps the individual model open to updating.
\end{itemize}

\section*{Self-Reinforcing Dynamics}

The interesting thing is that this can become self-reinforcing.

\begin{itemize}
    \item Better relationships allow better information exchange.
    \item Better information exchange allows better collective learning.
    \item Better collective learning can increase trust in the relationships that made it possible.
\end{itemize}

And the reverse is also possible:

\begin{itemize}
    \item Fear produces concealment.
    \item Concealment produces poorer information.
    \item Poorer information produces worse models.
    \item Worse models produce worse decisions.
    \item Worse decisions produce distrust.
    \item Distrust produces more fear and concealment.
\end{itemize}

So perhaps collective intelligence is not simply a property of how much information a group possesses.

Perhaps it is partly a property of whether the relationships between its members preserve the capacity for correction.

\section*{Decentralisation and Learning}

This also changes how I think about decentralisation.

If we do not know beforehand what the correct model of the future problem will be, we cannot simply distribute a correct solution from the centre.

Instead, we may need to distribute the \textit{capacity to learn}.

Not agreement on the answer.

Agreement on the practices that allow answers to be challenged, revised, and improved.

\section*{The Verbs of Learning}

This is why I have started thinking about the ``verbs of learning''.

Not what everyone should believe.

What everyone should be capable of doing.

\begin{itemize}
    \item Listening.
    \item Questioning.
    \item Explaining.
    \item Correcting.
    \item Updating.
    \item Admitting uncertainty.
    \item Seeking disconfirming information.
    \item Supporting someone else's learning.
    \item Being willing to change one's mind.
\end{itemize}

These behaviours do not automatically produce agreement. Nor do they guarantee that a collective will choose the right goal.

But they may create something more fundamental:

\begin{quote}
    A collective that remains capable of discovering that it is wrong.
\end{quote}

\section*{Reframing Decision Theory}

And perhaps that is the deeper problem I have been trying to describe.

Decision theory asks:

\begin{quote}
    What should we do?
\end{quote}

But under radical uncertainty, and especially when the decision is collective, perhaps we also need to ask:

\begin{quote}
    How do we remain capable of discovering that what we are doing is wrong?
\end{quote}

And if that capacity itself can be improved, then we arrive at a recursive problem:

\begin{itemize}
    \item We learn about the world.
    \item We learn how to learn about the world.
    \item And we learn together how to become better at learning together.
\end{itemize}

\section*{Conclusion}

I suspect that this is not a new idea in the literature. Many pieces of it almost certainly exist in different disciplines.

But I increasingly suspect that putting these pieces together points toward something important:

\begin{quote}
    The fundamental requirement for collective intelligence may not be agreement. It may be \textbf{collective corrigibility}.

    Not a collective that is permanently right.

    A collective that remains capable of becoming less wrong.
\end{quote}

\end{document}